\section{Contesto e Motivazione}
La tecnologia Blockchain ha ridefinito il concetto di fiducia digitale e decentralizzazione, permettendo lo scambio di valore e l'esecuzione di logiche complesse senza intermediari. Al centro di questa rivoluzione vi sono gli \textit{Smart Contract}, programmi immutabili che eseguono automaticamente clausole contrattuali al verificarsi di determinate condizioni. Tra le piattaforme di nuova generazione, \textbf{Algorand} si distingue per il suo protocollo di consenso \textit{Pure Proof-of-Stake} (PPoS), che garantisce scalabilità, sicurezza e finalità immediata delle transazioni.

Tuttavia, l'immutabilità della blockchain rappresenta un'arma a doppio taglio: una volta distribuito, il codice di uno smart contract non può essere modificato. Di conseguenza, eventuali vulnerabilità di sicurezza possono portare a perdite economiche ingenti e irreversibili. La sicurezza del codice diventa quindi un requisito non funzionale critico, specialmente in un ecosistema in cui gli asset gestiti hanno un valore monetario reale.

In ambito Algorand, lo sviluppo di smart contract avviene principalmente tramite l'Algorand Virtual Machine (AVM) e linguaggi di alto livello come \textbf{PyTeal} (un binding Python che compila in TEAL - \textit{Transaction Execution Approval Language}). La specificità e la complessità di questo stack tecnologico rendono l'audit di sicurezza particolarmente sfidante.

\section{Definizione del Problema}
Attualmente, la rilevazione di vulnerabilità negli smart contract si affida prevalentemente a due approcci, entrambi con limitazioni significative:

\begin{enumerate}
    \item \textbf{Strumenti di Analisi Statica}: Sebbene efficaci nell'individuare pattern noti e violazioni sintattiche, questi strumenti spesso mancano di "comprensione semantica". Tendono a generare un elevato numero di falsi positivi (segnalando errori dove non ci sono) o, peggio, falsi negativi (mancando vulnerabilità logiche complesse che dipendono dal contesto dell'applicazione).
    
    \item \textbf{Large Language Models (LLM)}: Modelli come GPT-5, Gemini o DeepSeek, dimostrano capacità di ragionamento eccezionali. Tuttavia, la loro performance degrada sensibilmente quando operano su linguaggi di dominio specifico e di nicchia come PyTeal. A differenza di Solidity (Ethereum), che è ampiamente rappresentato nei dataset di addestramento pubblici, PyTeal e TEAL sono meno diffusi. Di conseguenza, gli LLM generalisti tendono a "allucinare" sintassi inesistenti o a non comprendere le sottili implicazioni di sicurezza specifiche dell'AVM (es. gestione del \textit{rekeying} o utilizzo degli opcode).
\end{enumerate}

\section{Obiettivi e Contributo della Tesi}
L'obiettivo di questo lavoro di tesi è colmare il divario tra la potenza di ragionamento degli LLM moderni e la necessità di una conoscenza tecnica specifica per la sicurezza su Algorand. 

Al fine di superare tali limitazioni e ottimizzare le prestazioni degli attuali Large Language Models (LLM), in questo lavoro di tesi si propongono e si valutano due metodologie avanzate:

\begin{itemize}
    \item \textbf{Fine-tuning (Adattamento del Modello)}: Addestramento di un LLM open-source su un dataset curato di smart contract PyTeal e relative vulnerabilità note. Questo processo serve ad adattare il modello al linguaggio target, migliorandone la capacità di comprendere la sintassi e la struttura logica specifica di Algorand.
    
    \item \textbf{Retrieval-Augmented Generation (RAG)}: Integrazione del modello con una base di conoscenza esterna (Knowledge Base) contenente la documentazione aggiornata sulle vulnerabilità di Algorand. Il RAG permette al modello di accedere a informazioni di contesto mirate durante la fase di inferenza, riducendo le allucinazioni e garantendo che l'analisi sia conforme alle specifiche più recenti.
\end{itemize}

\section{Struttura dell'Elaborato}
Il presente documento è organizzato come segue:
\begin{itemize}
    \item Il \textbf{Capitolo 2} fornisce il background necessario sulla blockchain Algorand, il linguaggio PyTeal, le vulnerabilità comuni, una base teoriaca sugli LLM e RAG.
    \item Il \textbf{Capitolo 3} presenta lo stato dell'arte e ne descrive le problematiche e come tentiamo di migliorarle.
    \item Il \textbf{Capitolo 4} descrive la metodologia proposta, inclusa la preparazione del dataset, il processo di fine-tuning e l'implementazione della pipeline RAG.
    \item Il \textbf{Capitolo 5} presenta i risultati sperimentali.
    \item Il \textbf{Capitolo 6} trae le conclusioni e delinea possibili sviluppi futuri.
\end{itemize}